% !TeX root = ../../tesi.tex

\section{Il ruolo degli \emph{stage} nella strategia aziendale}
\label{sec:stage-strategia}

Halue ospita in media tre o quattro \emph{stage} all’anno. Gli stagisti provengono soprattutto
dall’Università di Padova, dai corsi di laurea triennale e magistrale, e alcuni degli \emph{stage}
che hanno preceduto il mio riguardavano già l’intelligenza artificiale. L’azienda attribuisce a
queste esperienze due scopi:

\begin{itemize}
	\item svolgere ricerca e sviluppo su temi che i progetti per i clienti non lasciano il tempo di
	esplorare;
	\item conoscere persone da inserire nel \emph{team}, valutando al termine di ogni percorso la
	possibilità di proseguire la collaborazione con lo stagista.
\end{itemize}

Il primo scopo spiega la forma che gli \emph{stage} assumono. Uno stagista lavora a tempo pieno su
un solo tema per circa due mesi, da 300 a 320 ore lavorative, con una continuità che i dipendenti
non hanno, perché l’azienda li ripartisce ogni mese tra più clienti
(\secref{sec:organizzazione-halue}). Halue affida quindi agli stagisti temi nuovi, su cui non ha
ancora un’esperienza consolidata, e un \emph{tutor} che li segue ogni giorno. Il secondo scopo
spiega perché lo stagista partecipa ai momenti di coordinamento del \emph{team} e ne adotta gli
strumenti, dato che l’azienda valuta il suo modo di lavorare oltre ai risultati che produce.

Ogni \emph{stage} si inserisce in un percorso che lo precede e lo segue. Il tema nasce da
un’esigenza interna oppure dall’interesse che può avere per un laureando in informatica, e nel mio
caso da entrambi, come illustro nella \secref{sec:obiettivi}. Al termine lo stagista presenta il
lavoro a tutto il \emph{team} e partecipa a una retrospettiva, la stessa pratica con cui Halue
chiude gli \emph{sprint} e i progetti per i clienti. Il lavoro può quindi avere due esiti:

\begin{itemize}
	\item l’azienda lo archivia come base di conoscenza sul tema;
	\item l’azienda lo rivede e ne riusa una parte nel codice e nei progetti di produzione.
\end{itemize}

La scelta dipende dal tipo di progetto e dalla qualità del prodotto.
